The two communications interfaces that can realistically transmit to the Onboard Computer (OBC) are I2C and SPI. Although USART was considered in \autoref{sec:v0.1-mcu}, it is not a high-speed serial protocol, so it will not be considered. The interfaces were considered using \href{https://resources.altium.com/p/spi-versus-i2c-how-choose-best-protocol-your-memory-chips}{this article}.
\nl
The OBC interface selection is driven by three primary features: speed, power consumption, and signal integrity. Faster interfaces will transfer data quicker, allowing the entire camera system to enter deep sleep quicker and save power. While the inherent power consumption of the interface is a consideration, the entire system has a higher consumption, so it is weighted considerably higher. Furthermore, the signal integrity of each interface is incredibly important, as a corrupted image is not desired.

\begin{indentedtable}
    \centering
    \begin{tabular}{|l|l|l|l|}
        \hline
        \textbf{Criteria} & \textbf{Weight} & \textbf{I2C} & \textbf{SPI} \\
        \hline
        Speed & 5 & 2 & 4 \\
        \hline
        Power Consumption & 1 & 3 & 4 \\
        \hline
        Signal Integrity & 3 & 3 & 5 \\
        \hline
        \hline
        \textbf{Total Score} &  & 22 & 39 \\
        \hline
    \end{tabular}
    \caption{C2Sv1 OBC interface protocol trade study.}
    \label{tab:v1-trade-obc}
\end{indentedtable}

From \autoref{tab:v1-trade-obc}, it is clear that SPI is the better protocol for transferring image data, which agrees with the \href{https://resources.altium.com/p/spi-versus-i2c-how-choose-best-protocol-your-memory-chips}{Altium article}. As the aim of the game is to transfer data in bulk with a tight window, power consumption from the interface alone becomes less of a concern as more power will be saved if the data is transferred quickly and the board can be put to sleep.